% HD-S37. GPT-6 Astra Pro, /root/sixfolds; inherited effort not exposed.
% 21 November 2026. Published classification with actual scalar/block maps.
\subsection{HD-S37: Gaussian sources of small codimension}
\label{sec:small-codimension-gaussian}

Use the actual primitive Jacobian factor and Gaussian corner specified
in HD-S36. Thus \(G\subset GL_m(K)\), \(K=\mathbb Q(i)\), is
finite and absolutely irreducible, its character field is exactly
\(K\), and the connected full-\(G\) sphere cover has a
signature-\((3,3)\) sixfold period family of rank nine.
The following assertions are proved here.

For \(m=9\), every reflection-free such coefficient group has
an actual \(K\)-defined monomial line system, and its ordinary
Hermitian determinant is \([-1]\). Together with HD-S36, all
coefficient ranks \(m\le9\) therefore have determinant
\([-1]\) or \([-3]\). For \(9\le m\le14\), a
reflection-free such coefficient representation is complex
imprimitive and has no nontrivial fixed tensor-product decomposition.
In particular the primitive degree-ten \(U_5(2)\) possibility
in the published small-codimension classification gives no new
Gaussian source. These statements make no assertion about the
remaining complex imprimitive groups of ranks ten through fourteen.

\paragraph{1. Numerical constraints and the exact published input.}
Write the nonidentity local monodromies as
\(g_1,\ldots,g_b\), with product one and generating \(G\),
and put \(c_j=\operatorname{rank}(g_j-1)\).
The actual parabolic Euler characteristic gives
\[
             \sum_{j=1}^b c_j=2m+6,\qquad b\ge12.             \tag{S37.1}
\]
The second inequality follows from the dimension \(b-3\) of the
branch configuration space and the asserted rank nine of the period
map. Reflection-free means \(c_j\ge2\). For \(m\le14\)
at least one \(c_j\) is two, since otherwise the left side of
(S37.1) is at least 36. For \(m=9\), (S37.1) gives
\[
                       b=12,\qquad c_j=2\quad(1\le j\le12).
                                                               \tag{S37.2}
\]

The published input is R.~M.~Guralnick--J.~Saxl,
\emph{Generation of finite almost simple groups by conjugates},
J.~Algebra \textbf{268} (2003), 519--571,
DOI \texttt{10.1016/S0021-8693(03)00182-0}.
Theorem~7.1, pp.~548--549, with its cross-characteristic
Lemma~7.4, p.~551, gives the following characteristic-zero
specialization: a finite complex primitive, tensor-indecomposable
representation of degree greater than eight containing a
bireflection is a scalar extension of a deleted permutation
representation, or has degree ten and normalizes the irreducible
degree-ten \(U_5(2)\) representation. The natural finite
classical alternatives in that theorem are in defining positive
characteristic. The shorter Theorem~7.A is printed on p.~521.
Theorem~8.4, pp.~557--558, will also be used: in degree nine an
irreducible imprimitive group generated by bireflections has a
normal abelian subgroup fixing nine lines, and its permutation
action on those lines is \(A_9\) or \(S_9\).
Only these specified specializations are used.

Here tensor indecomposability has the authors' stated meaning:
there is no invariant fixed tensor-product decomposition. A group
permuting equal tensor factors need not be tensor decomposable in
that meaning; such groups remain within their Theorem~7.1.
No classification of arbitrary finite Gaussian groups is being
silently substituted for these hypotheses.

\paragraph{2. Fixed tensor-product decompositions are excluded.}
Suppose that, over \(\mathbb C\), an actual intertwining
isomorphism \(T:V\to A\otimes B\) makes every element a
tensor product, where \(a=\dim A\le b'=\dim B\) and
\(ab'=m\). Uniqueness of a pure tensor up to reciprocal scalar
gives actual homomorphisms
\[
       G\longrightarrow PGL(A),\qquad G\longrightarrow PGL(B).
                                                               \tag{S37.3}
\]
Each factor of a finite-order tensor matrix is semisimple up to
scalar. If its \(A\)-factor is projectively nontrivial, every
eigenspace in that factor has dimension at most \(a-1\).
Summing over the eigenvalues of the \(B\)-factor shows that
every eigenspace of the tensor matrix has dimension at most
\((a-1)b'\). Thus its actual codimension is at least \(b'\).
Every nonidentity tensor matrix has codimension at least \(a\)
unless it is scalar, in which case its codimension is \(m\).

There must be at least three inertia images nontrivial in
\(PGL(A)\). With zero such images the full action on \(A\)
is scalar. With one, the product relation is impossible. With two,
the two images are inverse and generate a cyclic projective group;
choose a finite-order projective generator and a diagonalizable
lift. Every factor is a scalar times a power of this lift, so a
common eigenline in \(A\) gives an actual invariant subspace
under \(T^{-1}\) in \(V\). This contradicts irreducibility.

For \(9\le m\le14\), the factorizations with \(a\ge3\)
are \(3\cdot3\) and \(3\cdot4\); their total codimension
is at least \(12a>2m+6\). The remaining factorizations are
\(2\cdot b'\), with \(b'=5,6,7\). At least three inertia
elements have codimension at least \(b'\), and all others at
least two. Consequently
\[
        \sum c_j\ge 3b'+2(12-3)=3b'+18>4b'+6=2m+6.
                                                               \tag{S37.4}
\]
All fixed tensor decompositions in the stated ranks are excluded.

\paragraph{3. Actual scalar maps for deleted permutations.}
In the deleted permutation alternative set \(N=m+1\), and
let \(P_\pi\) denote the permutation \(\pi\in S_N\)
restricted to the augmentation space
\(V_0=\{(z_1,\ldots,z_N):\sum z_i=0\}\).
In characteristic zero the constant line splits off the permutation
module, so the deleted module here has dimension \(N-1\).
The classified representation supplies a complex intertwiner with
\[
                  T\rho(g)T^{-1}=a(g)P_{\pi(g)}.             \tag{S37.5}
\]
These are literal group maps: a permutation acting as a scalar on
\(V_0\), for \(N\ge3\), is the identity (apply it to the
vectors \(e_j-e_k\)). Hence the pair in (S37.5) is unique,
\(\pi(gh)=\pi(g)\pi(h)\), and \(a(gh)=a(g)a(h)\).
Finiteness of \(G\) makes every \(a(g)\) a root of unity.
These maps and their eigenvalue computations are on the actual
representation; no rational descent of this complex intertwiner
is needed for the following contradiction.

If all inertia phases are real, they are \(\pm1\), and the
full generated character is real. Since its values lie in
\(K\), its character field would then be \(\mathbb Q\).
There are therefore at least two nonreal inertia phases: exactly
one cannot have product one with the remaining real phases.
For a nonreal phase, its inverse has order at least three.
A cycle of a permutation contains this eigenvalue once or zero
times, and only if that order divides its length. Its multiplicity
on \(V_0\) is therefore at most \(\lfloor N/3\rfloor\).
The corresponding actual inertia codimension is at least
\(m-\lfloor N/3\rfloor\). Equation (S37.1) now gives
\[
 2m+6\ge2\bigl(m-\lfloor N/3\rfloor\bigr)+2(b-2),
 \qquad b\le5+\lfloor N/3\rfloor\le10<12                  \tag{S37.6}
\]
for \(9\le m\le14\). Thus this classified primitive
alternative cannot supply the stated Gaussian family.

\paragraph{4. The degree-ten exceptional module and its normalizers.}
The \(U_5(2)\) ordinary character table has a unique irreducible
character of degree ten, \(\chi_2\), which is rational-valued
and has Frobenius--Schur indicator \(-1\).
The primary ATLAS representation entry
\texttt{U52G1-Ar10B0} supplies its actual Gaussian integral
generators \(x,y\) and a nondegenerate alternating tensor
\(B_1\) with determinant one. The entry is
\url{https://brauer.maths.qmul.ac.uk/Atlas/v3/matrep/U52G1-Ar10B0};
the inert matrix data are
\url{https://brauer.maths.qmul.ac.uk/Atlas/clas/U52/mag/U52G1-Ar10B0.M}.
The ordinary character-table input is also identified exactly by
\texttt{CharacterTable("U5(2)")}, irreducible row 2, in the
ATLAS-derived Character Table Library. Its complete degree-ten
row and indicator were checked; no claim of reconstructing that
entire character table is made.

The exact accompanying certificate parses these matrices as data,
and verifies
\[
 xB_1x^t=yB_1y^t=B_1,\qquad \det B_1=1,\qquad
 J=B_1^{-1},\quad x^tJx=y^tJy=J.                              \tag{S37.7}
\]
It also verifies a positive invariant Hermitian tensor \(B_2\),
whose leading principal minors are
\[
        4,14,24,68,72,136,96,112,64,32.
\]
Gaussian integrality and this positive form prove finiteness of
the generated matrix group. The seven defining relators in the
primary ATLAS presentation of \(U_5(2)\) are checked literally.
The commutant equations have rank 99 modulo \(5\), under
\(i\mapsto2\), and have the scalar matrices in their kernel;
they therefore have complex kernel dimension one. Finiteness
then proves absolute irreducibility. This separately verifies the
existence and uniqueness up to scalar of the invariant alternating
form on the actual degree-ten model used by the classification.

Every normalizing matrix \(g\) satisfies
\(g^tJg=\lambda(g)J\): the left side is again invariant
under the irreducible normal subgroup, and the invariant-form
space is one-dimensional. The multiplier \(\lambda\) is a
literal character. In degree ten, (S37.1), \(b\ge12\), and
\(c_j\ge2\) give \(c_j\le26-22=4\).
If \(\lambda(g_j)\ne1\), its fixed subspace is isotropic
for \(J\), since for fixed vectors \(u,v\)
\[
             J(u,v)=J(g_ju,g_jv)=\lambda(g_j)J(u,v).
\]
That subspace has dimension at least six, whereas an isotropic
subspace for this ten-dimensional nondegenerate alternating form
has dimension at most five. Hence every inertia multiplier is one.
Their generated group preserves \(J\), so its representation is
self-dual and its finite-order eigenvalues occur with reciprocal
multiplicity. Its character is real, again contradicting the
exact character field \(\mathbb Q(i)\).
This treats every normalizer allowed by the published exceptional
case, without identifying an outer extension with a scalar extension.

Paragraphs 2--4 prove that all reflection-free coefficient groups
in ranks nine through fourteen under consideration are complex
imprimitive, and exclude the exceptional degree-ten candidate.

\paragraph{5. Descent of the rank-nine complex line system.}
Now put \(m=9\). The group is generated by the twelve actual
bireflections in (S37.2). Theorem~8.4 gives nine complex lines
\(L_1,\ldots,L_9\), their pointwise stabilizer
\(T_0\triangleleft G\), and the actual quotient
\[
       P=G/T_0=A_9\ \text{or}\ S_9
       \quad\text{on }\{L_1,\ldots,L_9\}.                  \tag{S37.8}
\]
The diagonal group \(T_0\) is abelian. We first prove that it
cannot consist only of scalars, rather than assuming the lines
are already defined over \(K\).

For a monomial cycle of length \(\ell\) with holonomy \(u\),
the fixed dimension is one for \(u=1\), and zero otherwise.
Consequently the permutation of a bireflection is the identity,
a transposition, a 3-cycle, or two disjoint transpositions.
If \(T_0\) is scalar, the identity permutation is impossible.
For a transposition permutation, the square of the matrix is in
\(T_0\); at least six unaffected lines have eigenvalue one,
so that square is the identity. The two-cycle holonomy is then
one, and exactly one of the seven fixed lines has eigenvalue
\(-1\), the other six having eigenvalue one. Its trace is 5.
The centralizer of the transposition in \(S_9\) acts
transitively on those seven lines. Lift such permutations to
\(G\). Each lift conjugates the bireflection to a scalar
multiple of itself; equality of the nonzero traces forces that
scalar to be one. This would make the seven fixed-line
eigenvalues equal, a contradiction. The \(A_9\) quotient has
no transposition in the first place.

The remaining inertia permutations are 3-cycles or double
transpositions; all their cycle holonomies and all unaffected
eigenvalues are one. Thus they have determinant one and even
permutation. They generate a group with quotient \(A_9\).
Its scalar subgroup lies in \(K^*\), so is contained in
\(\mu_4\); determinant one also gives \(\zeta^9=1\),
forcing \(\zeta=1\). Hence \(G\simeq A_9\).
The stabilizer of one line is \(A_8\), which is perfect:
for example the commutator of \((12)(45)\) and
\((23)(45)\) is a 3-cycle, and its conjugates generate \(A_8\).
Its line character is therefore trivial. The actual coset map
\[
    \mathbb C[G]\otimes_{\mathbb C[A_8]}\mathbb C
       \longrightarrow V,\qquad g\otimes1\longmapsto\rho(g)v
                                                               \tag{S37.9}
\]
is an isomorphism, for nonzero \(v\) on that line: it sends
the nine cosets to nine independent lines. The sum of the coset
basis is a nonzero invariant vector. This contradicts
irreducibility, completing the exclusion of scalar \(T_0\).

Let \(\chi_j:T_0\to\mathbb C^*\) be the character on
\(L_j\). Equality of two such characters is an invariant
equivalence relation under the 2-transitive group \(P\).
Since \(T_0\) is not scalar, all nine characters are distinct.
The lines are therefore exactly its one-dimensional common
eigenspaces. All \(T_0\)-matrices have entries in \(K\),
and the character values are roots of unity. A Galois automorphism
of their cyclotomic splitting field over \(K\) consequently
permutes these nine characters and lines. This permutation
commutes with \(P\), because every \(\rho(g)\) is fixed
by that automorphism. The centralizer of the natural \(A_9\)
or \(S_9\) in \(S_9\) is trivial: its point stabilizer
fixes only that point, and transitivity then fixes every point.
Thus every \(\chi_j\) is \(K\)-valued, and
\[
        p_j=\frac1{|T_0|}\sum_{t\in T_0}\chi_j(t)^{-1}\rho(t)
            \in M_9(K),\qquad L_j(K)=\operatorname{im}p_j     \tag{S37.10}
\]
has rank one. These are the promised actual \(K\)-lines.

\paragraph{6. The coefficient and Jacobian morphisms to HD-S34.}
Choose \(0\ne v\in L_1(K)\), let \(H\le G\) stabilize
that line, and let \(\chi:H\to K^*\) be its action.
Because \(H\) is finite, \(\chi(H)\subset\mu_4\).
For left coset representatives \(t_j\), the map
\[
     K^9\longrightarrow V,\qquad e_j\longmapsto\rho(t_j)v
                                                               \tag{S37.11}
\]
is an actual \(K\)-intertwiner from \(\operatorname{Ind}_H^G\chi\).
If \(gt_j=t_kh\), its matrix has entry \(\chi(h)\)
from the \(j\)-th to the \(k\)-th line. The induced trace
formula and the exact original character field force
\(\chi(H)=\mu_4\). For any positive invariant Hermitian
form, the distinct \(T_0\)-character lines are orthogonal,
and the nine chosen vectors have the same squared norm
\(q=h(v,v)>0\). Thus (S37.11) identifies the coefficient
forms with the literal form \(qI_9\).

This is precisely the \(K\)-defined line-system hypothesis of
HD-S34, now proved from the complex classification. To retain
the actual curve maps, put
\[
      Y=\widetilde C/H,\qquad
      X=\widetilde C/\ker\chi,\qquad X\longrightarrow Y.
                                                               \tag{S37.12}
\]
The latter is the actual cyclic-four quotient of the original
cover. The primitive line idempotent on \(J\widetilde C\)
equals the lift of the \(\chi+\bar\chi\) idempotent on
\(JX\), after the \(\ker\chi\)-average, as proved by
the same literal coefficient formula in HD-S34. Pullback
\(JX\to J\widetilde C\) is a polarized isogeny on that
factor with multiplier \(|\ker\chi|\); the inverse is its
norm divided by \(|\ker\chi|\) on rational homology.
The common noncentral Gaussian corner is the \(h_0\)-action
of HD-S34 for \(\chi(h_0)=i\). Hence these maps identify
the specified \(K\)-actions, not just the unpolarized complex
factors. The multiplier on the six-dimensional Hermitian
determinant is \(|\ker\chi|^6=N(|\ker\chi|^3)\).
HD-S34's cyclic-cover cup computation therefore gives ordinary
determinant \([-1]\), and HD-S03a identifies the resulting
signature-\((3,3)\) rational space with the split space.
All permitted coefficient-line and fixed-polarization transports
retain exactly the maps of HD-S17 and HD-S36.

\paragraph{7. Evidence and full author audit.}
The exact data and calculations are
\path{work/sixfolds/small_codimension/small_codimension_certificate.py}
and its JSON/log, with the original inert
\texttt{U52G1-Ar10B0.M} data.
The separate GAP probe records the exact unique degree-ten
character and Frobenius--Schur indicator; it is a published-table
check, not an independent character-table construction.
The complete author proof check covers (S37.1)--(S37.12), including
the two distinct scalar arguments, the actual cyclic projective
factor map in the tensor case, the scalar-diagonal impossibility,
Galois descent of each projector, and the polarized quotient
morphisms. The remaining complex imprimitive coefficient groups
of ranks ten through fourteen are left as explicit open source
gates, with no dominance or determinant assertion made for them.
